\documentclass[11pt, a4paper]{article}

% A formal letter. Replace the addresses, date, subject and text below; the
% layout follows the usual block style with everything aligned left except
% the sender's address.

% Fonts and encoding: pdfLaTeX needs these; XeLaTeX and Tectonic do not.
\usepackage{iftex}
\ifPDFTeX
  \usepackage[T1]{fontenc}
  \usepackage{lmodern}
\fi
\usepackage{microtype}
\usepackage[margin=1in]{geometry}
\usepackage{xcolor}
\usepackage[colorlinks=true, urlcolor=blue!60!black]{hyperref}

\pagestyle{empty}
\setlength{\parindent}{0pt}
\setlength{\parskip}{0.8em}

\begin{document}

\begin{flushright}
Your Name \\
Street Address \\
City, Postcode \\
Country \\
\href{mailto:you@example.org}{you@example.org}
\end{flushright}

\today

Recipient's Name \\
Job Title \\
Organisation \\
Street Address \\
City, Postcode

\textbf{Subject: A short line saying what the letter is about}

Dear Recipient's Name,

Open with the reason you are writing, in one or two sentences, so the
reader knows at once what the letter is about and what you are asking for.

Use the middle paragraphs to give the details: the background, the facts
that support your request, and any dates or reference numbers the reader
will need. Keep each paragraph to a single point.

Close by saying what you would like to happen next and how the reader can
reach you. Thank them for their time.

Yours sincerely,

\vspace{2.5em}

Your Name \\
Your Job Title

\textbf{Enclosures:} list any documents you are sending with the letter.

\end{document}
